\documentclass{ximera}

\input{../../preamble.tex}

\title{Chain Rule}

\begin{document}

\begin{abstract}
behavior
\end{abstract}
\maketitle




The Chain Rule can tell us about behavior.


What is the behavior of 

\[ H(x) = 3 \sin(3 - 2x)  - 4 \]


We can view this as a composition: 


\[ H(x) = 3 \sin(3 - 2x)  - 4 =   (Out \circ \sin \circ In)(x)\]

where


\begin{itemize}
\item $Out(t) = 3 \, t - 4$ is an increasing linear function
\item $\sin(\theta)$ is the core sine function
\item $In(y) = 3 - 2 \, y$ is a decreasing linear function
\end{itemize}



The core sine function has its behavior:



\begin{itemize}
\item $\sin(\theta)$ increases on $\left( 0, \frac{\pi}{2} \right)$
\item $\sin(\theta)$ decreases on $\left( \frac{\pi}{2}, \pi \right)$
\item $\sin(\theta)$ decreases on $\left( \pi, \frac{3\pi}{2} \right)$
\item $\sin(\theta)$ increases on $\left( \frac{3\pi}{2}, 2\pi \right)$
\end{itemize}



We need to know where $In(y)$ equals these endpoints.



$\blacktriangleright 3 - 2y = 0$ at $\frac{3}{2}$

$\blacktriangleright 3 - 2y = \frac{\pi}{2}$ at $-\frac{1}{2}(\frac{\pi}{2} - 3) = \frac{6-\pi}{4} = \frac{3}{2} - \frac{\pi}{4}$


\begin{center}
The first quadrant of sine corresponds to $\left( \frac{3}{2} - \frac{\pi}{4}, \frac{3}{2} \right)$.

\[ Out \circ \sin \circ In = inc \circ inc \circ dec = dec \]
\end{center}




$\blacktriangleright 3 - 2y = \pi$ at $-\frac{1}{2} (\pi - 3) = \frac{3-\pi}{2} = \frac{3}{2} - \frac{\pi}{2}$


\begin{center}
The second quadrant of sine corresponds to $\left( \frac{3}{2} - \frac{\pi}{2}, \frac{3}{2} - \frac{\pi}{4} \right)$.

\[ Out \circ \sin \circ In = inc \circ dec \circ dec = inc \]
\end{center}




$\blacktriangleright 3 - 2y = \frac{3\pi}{2}$ at $-\frac{1}{2}(\frac{3\pi}{2} - 3) = \frac{6-3\pi}{4} = \frac{3}{2} - \frac{3\pi}{4}$


\begin{center}
The third quadrant of sine corresponds to $\left( \frac{3}{2} - \frac{3\pi}{2}, \frac{3}{4} - \frac{\pi}{2} \right)$.

\[ Out \circ \sin \circ In = inc \circ dec \circ dec = inc \]
\end{center}






$\blacktriangleright 3 - 2y = 2\pi$ at $-\frac{1}{2} (2\pi - 3) = \frac{3-2\pi}{2}  \frac{3}{2} - \pi$


\begin{center}
The fourth quadrant of sine corresponds to $\left( \frac{3}{2} - \pi, \frac{3}{2} - \frac{3\pi}{4} \right)$.

\[ Out \circ \sin \circ In = inc \circ inc \circ dec = dec \]
\end{center}








The numbers are coming out in descending order because of the negative leading coefficient.

The period is $\frac{3}{2} - \frac{3-2\pi}{2} = \pi$




\textbf{\textcolor{blue!55!black}{On $\left( \frac{3}{2} - \pi, \frac{3}{2} - \frac{3\pi}{4} \right)$}} $H$ is decreasing.


\textbf{\textcolor{blue!55!black}{On $\left( \frac{3}{2} - \frac{3\pi}{4}, \frac{3}{2} - \frac{\pi}{2} \right)$}} $H$ is increasing.


\textbf{\textcolor{blue!55!black}{On $\left( \frac{3}{2} - \frac{\pi}{2}, \frac{3}{2} - \frac{\pi}{4} \right)$}} $H$ is increasing.


\textbf{\textcolor{blue!55!black}{On $\left( \frac{3}{2} - \frac{\pi}{4}, \frac{3}{2} \right)$}} $H$ is decreasing.


\begin{image}
\begin{tikzpicture}
  \begin{axis}[
            domain=-10:10, ymax=10, xmax=10, ymin=-10, xmin=-10,
            axis lines =center, xlabel=$x$, ylabel={$y=H(x)$}, grid = major, grid style={dashed},
            ytick={-10,-8,-6,-4,-2,2,4,6,8,10},
            xtick={-10,-8,-6,-4,-2,2,4,6,8,10},
            yticklabels={$-10$,$-8$,$-6$,$-4$,$-2$,$2$,$4$,$6$,$8$,$10$}, 
            xticklabels={$-10$,$-8$,$-6$,$-4$,$-2$,$2$,$4$,$6$,$8$,$10$},
            ticklabel style={font=\scriptsize},
            every axis y label/.style={at=(current axis.above origin),anchor=south},
            every axis x label/.style={at=(current axis.right of origin),anchor=west},
            axis on top
          ]
          

			   \addplot [line width=2, penColor, smooth,samples=200,domain=(-6:7)] {3*sin(deg(3-2*x)) - 4};

          \addplot [line width=2, penColor2, smooth,samples=200,domain=(-1.64:1.5)] {3*sin(deg(3-2*x)) - 4};

  \end{axis}
\end{tikzpicture}
\end{image}




Just Checking...

\textbf{\textcolor{blue!55!black}{On $(-1.64, -0.86)$}} $H$ is decreasing.


\textbf{\textcolor{blue!55!black}{On $(-0.86, -0.07)$}} $H$ is increasing.


\textbf{\textcolor{blue!55!black}{On $(-0.07, 0.71)$}} $H$ is increasing.


\textbf{\textcolor{blue!55!black}{On $(0.71, 1.5)$}} $H$ is decreasing.














































\begin{center}
\textbf{\textcolor{green!50!black}{ooooo-=-=-=-ooOoo-=-=-=-ooooo}} \\

more examples can be found by following this link\\ \link[More Examples of Easy Angles]{https://ximera.osu.edu/csccmathematics/precalculus2/precalculus2/easyAngles/examples/exampleList}

\end{center}





\end{document}

